\documentclass[11pt]{article}
\usepackage[utf8]{inputenc}
\usepackage{amsmath,amsthm,booktabs}
\newtheorem{proposition}{Proposición}
\input{../output/macros.tex}

\title{Restricción de crédito y transmisión monetaria a la actividad industrial}
\date{Borrador}

\begin{document}
\maketitle

\section{Pregunta}
¿Reduce un endurecimiento monetario la producción industrial de EE.~UU. más cuando el crédito corporativo está caro?
Estimamos proyecciones locales mensuales (1990--2019) de la producción industrial ante los choques de política monetaria
de Jarociński y Karadi, con una interacción con el diferencial Baa del mes anterior.

\section{Teoría}
Una empresa con tecnología $A k^\alpha$ afronta un coste del capital $c = r(1 + m s) + \delta$, donde $s \geq 0$ es el
diferencial de crédito y $m > 0$ la sensibilidad de la prima de financiación externa, que escala con el coste de la
deuda. El capital óptimo es $k^* = (\alpha A / c)^{1/(1-\alpha)}$.

\begin{proposition}
La semielasticidad del capital deseado al tipo de interés aumenta con el diferencial de crédito:
\[
  \frac{\partial^2 \log k^*}{\partial r \, \partial s} = \frac{m \delta}{(1-\alpha)\, c^2} > 0 .
\]
\end{proposition}

La proposición motiva el signo esperado de la interacción; no deriva la proyección local, que es una forma reducida.
El signo depende de que la prima sea proporcional al tipo: con una prima aditiva se invierte.

\section{Resultados}
La especificación prerregistrada es la proyección local a $h = 12$ meses con \lags{} retardos de la producción, la
inflación, el paro, el diferencial y el propio choque, y errores robustos con residuos HC3 ($N = \nobsHtwelve$). El
estado entra centrado en su media muestral, así que $b_h$ es la respuesta con el diferencial en esa media. Ante un
endurecimiento de una desviación típica, la producción a un año varía $\bShockHtwelve$ puntos logarítmicos
(e.e.~$\seShockHtwelve$). La interacción es $c_{12} = \bInterHtwelve$ (e.e.~$\seInterHtwelve$; intervalo bilateral al
\ciLevel{}~\% $[\loInterHtwelve, \hiInterHtwelve]$; $p$ unilateral $\pInterHtwelve$): no se confirma H1.

\begin{table}[h]
\centering
\begin{tabular}{lccccc}
\toprule
$h$ (meses) & 0 & 6 & 12 & 18 & 24 \\
\midrule
$b_h$ (choque) & $\bShockHzero$ & $\bShockHsix$ & $\bShockHtwelve$ & $\bShockHeighteen$ & $\bShockHtwentyfour$ \\
 & ($\seShockHzero$) & ($\seShockHsix$) & ($\seShockHtwelve$) & ($\seShockHeighteen$) & ($\seShockHtwentyfour$) \\
$c_h$ (choque $\times$ crédito) & $\bInterHzero$ & $\bInterHsix$ & $\bInterHtwelve$ & $\bInterHeighteen$ & $\bInterHtwentyfour$ \\
 & ($\seInterHzero$) & ($\seInterHsix$) & ($\seInterHtwelve$) & ($\seInterHeighteen$) & ($\seInterHtwentyfour$) \\
\bottomrule
\end{tabular}
\caption{Proyecciones locales de la producción industrial acumulada. Errores HC3 entre paréntesis.}
\end{table}

Un Monte Carlo con el mismo estimador, $T$ y un proceso calibrado a estos datos (choques remuestreados, con su masa en
cero y sus colas; \texttt{output/MC\_REPORT.md}) da una cobertura dentro de la banda prerregistrada, así que usamos estos
intervalos. El mismo Monte Carlo muestra que la potencia frente a una interacción del tamaño simulado es baja: el
resultado nulo no descarta efectos de ese orden.

\section{Conclusión}
Con choques de alta frecuencia y datos mensuales, no encontramos evidencia de que el diferencial de crédito amplifique la
transmisión de la política monetaria a la producción industrial, aunque la muestra tiene poca potencia para detectarlo.

\end{document}
